\section{DISCUSSION}\label{discussion}

\subsection{Principal findings}\label{principal-findings}

The headline result is that \textbf{fabrication is bounded in modern
frontier LLMs and stratifies cleanly by model tier and prompting
mode}, while abstention is the more consequential failure mode
for some models. The 6-way taxonomy is what makes the second
pattern visible. Across the 24 LLM configurations, hallucination
ranges from 0\% to 5.6\% for flagship models (Opus 4.7, GPT-5.5,
Gemini 2.5 Pro) and from 0\% to 12.8\% for sub-flagship models, with
the upper bound from GPT-4o-mini zero-shot D2. Hallucination
behaves as a property of the prompting setup rather than a fixed
property of the model. The shift from zero-shot to constrained
prompting eliminates fabrication entirely for every Anthropic and
OpenAI model that shows any nonzero zero-shot hallucination, and
within-model paired McNemar tests confirm significance at n=125.
The string-matching baselines and the frozen sentence-transformer
baseline (45 to 69\% top-1) sit well below every constrained-mode
LLM on the fabrication-versus-accuracy frontier, confirming that
constrained frontier LLMs offer real gains over lexical and
embedding-only approaches at this cohort scale. Cost-per-correct
collapses the deployment choice to a small frontier: GPT-4o-mini
constrained at \$0.0003 and Claude Haiku 4.5 constrained at
\$0.0044 dominate every more expensive configuration tested.

\subsection{Comparison with prior work}\label{comparison-with-prior-work}

The closest direct precedent is Soroush et al.~(2024)
\autocite{Soroush2024}, who benchmarked four 2023--2024 vintage
LLMs on direct ICD code querying. Our matrix differs from theirs in
three ways relevant to deployment. The first is model generation.
Their headline numbers are GPT-4 (46\% exact match on ICD-9, 34\%
on ICD-10); our matrix includes Opus 4.7, GPT-5.5, and Gemini 3
Flash Preview on the same task type. Every Anthropic and OpenAI
configuration we evaluated exceeds 80\% top-1 even under the D4
abbreviation stress, which we attribute to cohort-scope
restriction, generational LLM improvements, and structured-output
prompting. The second is per-prediction cost. Soroush et al.\ and
the broader hallucination-mitigation literature
\autocite{DasBaksi2024,Sarvari2025,Motzfeldt2025,HybridCode2025} do
not report per-prediction cost; we do, and cost-per-correct
collapses the choice to a small Pareto frontier (GPT-4o-mini
constrained at \$0.0003 and Claude Haiku 4.5 constrained at
\$0.0044). The third is the outcome taxonomy: hallucination is an
explicit, mutually-exclusive bucket, abstention is split out via
the 6-way taxonomy, and the within-model zero-shot vs.\ constrained
vs.\ RAG contrast is reported as a controlled ablation that
separates ``wandering off the menu'' from genuine ignorance.
Neuro-symbolic alternatives \autocite{Motzfeldt2025,HybridCode2025}
report hallucination control through architectural constraints;
our 0\% Anthropic constrained-mode rate sits below their reported
neural-baseline band, suggesting that modern frontier LLMs under
structured-output constraints have closed much of the fabrication
gap those systems were designed to fix. Full comparison is in
Supplementary §S2.

\subsection{Implications for clinical deployment}\label{implications-for-clinical-deployment}

The 6-way taxonomy reframes the deployment question. Top-1
accuracy framing treats every wrong answer equally, but the cost
of a hallucinated code (downstream systems silently mishandle it,
propagating into billing, research, and quality reporting) is
qualitatively different from a near-miss within the correct ICD
chapter (corrected by a reviewer in seconds), which is in turn
different from an abstention (re-coding from scratch, but
emitting no spurious artifact).

The deployment pattern supported by these results is therefore
\textbf{LLM-augmented coding with terminologist or
clinical-informaticist review}, not autonomous LLM coding. Two
findings concretely support this. For constrained-mode Anthropic
and OpenAI configurations (0\% hallucination, \textasciitilde5\% category-match,
\textasciitilde0\% abstention), the QA fraction is dominated by fast
category-match adjudication, a clinically straightforward task an
experienced coder can resolve in seconds, which beats human-only
coding across all sensitivity-sweep parameter ranges considered
(Supplementary §S5). The top-1 vs top-5 lift table (§4.4) also
shows that trained-classifier and sub-frontier zero-shot LLM
workflows (top-1 in the 45 to 75\% range) benefit substantially
from surfacing top-5 candidates to a reviewer, lifting net
accuracy 20 to 38 percentage points. Constrained-mode frontier
LLMs already saturate top-1 (lift $\leq$6pp) and shift the reviewer's
role from candidate selection to category-match adjudication.
Either architecture keeps a clinical adjudicator in the decision
loop where accuracy matters most.

\subsection{Limitations}\label{limitations}

The v1 headline matrix evaluates n=125 unique Synthea Conditions.
Per-cell Wilson 95\% CIs on a 10\% rate are correspondingly ±5pp,
and v2 replication at n=2000 will tighten intervals by
approximately 2×. Gemini 2.5 Pro showed an 88\% no\_prediction
rate on zero-shot D4 with most empty rows showing zero output
tokens, consistent with reasoning-token budget exhaustion under
our wrapper's \texttt{max\_output\_tokens=1024} setting. Whether this is
a fixable wrapper interaction or a genuine model limitation is
the subject of a follow-up investigation blocked by Tier-1 daily
quota at v1 submission. Gemini 3.1 Pro Preview was excluded for
the same quota reason. The benchmark covers English ICD-10-CM
only; extension to LOINC and RxNorm via the same taxonomy is
planned for v2.

\subsection{Future directions}\label{future-directions}

Two extensions follow directly. The first is v2 vocabulary
expansion to LOINC and RxNorm; the 6-way taxonomy and eval runner
already support arbitrary FHIR resource types. The second is
rationale evaluation following Li et al.~\autocite{Li2025}, in
which each LLM is asked for a per-prediction rationale to support
faithfulness and plausibility evaluation alongside the
outcome-bucket framework. A v2 cohort with broader vocabulary and
more realistic clinical text (e.g., discharge summaries, ED notes)
is also needed to establish the true effect size of fabrication
risk on real clinical inputs.
